\documentclass[12pt]{article}
\usepackage[T1]{fontenc}
\usepackage[utf8]{inputenc}
\usepackage{lmodern}
\usepackage{textcomp}
\usepackage{enumitem}
\usepackage{geometry}
\geometry{margin=1in}
\begin{document}
\begin{center}
Answer Key - Psychology - Undergraduate Level Assessment 19 \
\small (Answers are ordered exactly as in the question paper.)
\end{center}

\section*{Multiple Choice}
\begin{enumerate}[label=\arabic*.]
\item \textbf{Answer:} B
\\ \textit{Options:} A) 5; B) 8; C) 11; D) 12; E) 4
\\ \textit{Note:} The median is the middle value of a data set when arranged in ascending order, and for the numbers 3, 5, 8, 9, and 12, the median is 8, as it is the third number in the ordered list.
\item \textbf{Answer:} E
\\ \textit{Options:} A) "I don't want to waste gas."; B) "I'm not in a rush."; C) "I want to impress my friends."; D) "I find it relaxing to drive slowly."; E) "It's the law."
\\ \textit{Note:} The correct answer illustrates Kohlberg's conventional level of morality because it reflects adherence to societal rules and laws, indicating that the 17-year-old values conformity and the maintenance of social order in his decision-making.
\item \textbf{Answer:} A
\\ \textit{Options:} A) has no relation with the individual's personal life; B) tends to decrease steadily as tenure on a job increases; C) is always high in jobs that require high levels of education; D) is solely dependent on the job's financial compensation; E) is only achieved if an individual is in a management position
\\ \textit{Note:} The correct answer reflects the understanding that job satisfaction is primarily influenced by factors related to the work environment and organizational dynamics, rather than personal life circumstances, which may not directly impact one's satisfaction with their job.
\item \textbf{Answer:} B
\\ \textit{Options:} A) that language is a uniquely human skill.; B) the importance of physical contact to development.; C) that primates have a hierarchy based on physical strength.; D) that monkeys have a preference for inanimate objects.; E) that primates can learn complex tasks.
\\ \textit{Note:} Harlow's monkey experiment demonstrated that infant monkeys preferred to cling to a soft, comforting surrogate mother over a wire mother that provided food, highlighting the crucial role of physical contact and emotional comfort in healthy psychological development.
\item \textbf{Answer:} A
\\ \textit{Options:} A) constant changes in levels of arousal.; B) low levels of arousal.; C) arousal levels that are consistently low.; D) high levels of arousal.; E) complete absence of arousal.
\\ \textit{Note:} The correct answer is based on the Yerkes-Dodson law, which posits that optimal learning and performance occur at moderate levels of arousal, and that fluctuations in arousal can enhance engagement and adaptability to varying tasks and challenges.
\end{enumerate}

\section*{True / False}
\begin{enumerate}[label=\arabic*.]
\setcounter{enumi}{5}
\item \textbf{Answer:} True
\\ \textit{Options:} A) True; B) False
\\ \textit{Note:} The statement is true because a phoneme is the smallest unit of sound in spoken language, and the 'c' in 'cat' represents a specific sound that distinguishes it from other words.
\item \textbf{Answer:} False
\\ \textit{Options:} A) True; B) False
\\ \textit{Note:} The variance of the sample observations 2, 5, 7, 9, 12 is calculated to be 10.8, not 12.4, indicating that the statement is false.
\item \textbf{Answer:} True
\\ \textit{Options:} A) True; B) False
\\ \textit{Note:} In Bandura's social learning theory, self-reinforcement allows individuals to control and sustain their behavior by providing internal rewards for their actions, independent of external reinforcement.
\item \textbf{Answer:} True
\\ \textit{Options:} A) True; B) False
\\ \textit{Note:} Job productivity is a significant predictor of employees' long-term retention because higher productivity often leads to increased job satisfaction, recognition, and opportunities for advancement, which foster a stronger commitment to the organization.
\item \textbf{Answer:} False
\\ \textit{Options:} A) True; B) False
\\ \textit{Note:} Object permanence is actually a concept associated with the sensorimotor stage of Piaget's theory, where infants learn that objects continue to exist even when they cannot be seen, while the concrete operational stage involves the development of logical thinking and understanding of conservation.
\end{enumerate}

\section*{Long Form Response}
\begin{enumerate}[label=\arabic*.]
\setcounter{enumi}{10}
\item \textbf{Answer:} Catharsis in psychology refers to the process of releasing pent-up aggressive energy, allowing individuals to alleviate emotional tension. Theorists suggest that aggressive impulses can be effectively managed by redirecting this energy into more productive or neutral outlets, such as physical exercise, creative activities, or engaging in problem- solving tasks. By channeling aggression into these constructive avenues, individuals can mitigate the potential for harmful behaviors while experiencing emotional relief. This approach posits that rather than suppressing aggressive feelings, acknowledging and transforming them can lead to healthier emotional regulation and well-being.
\\ \textit{Note:} correct because it accurately describes catharsis as a process of releasing aggressive energy and explains how redirecting this energy into constructive activities can help manage aggressive impulses and reduce the likelihood of harmful behaviors.
\item \textbf{Answer:} Carl Rogers' theory of psychological development emphasizes the significant impact of external evaluations on an individual's self-concept, particularly during childhood. He posits that when children receive conditional positive regard-where their worth is evaluated based on specific behaviors or achievements-they may internalize these judgments, leading to a disconnect between their true self and the self they present to the world. This incongruence results in psychological maladjustment, as the individual struggles to reconcile their genuine experiences with the expectations imposed by others. Consequently, they may develop low self-esteem, anxiety, or other maladaptive behaviors, illustrating how a lack of unconditional positive regard can hinder personal growth and well-being. Ultimately, Rogers highlights the importance of fostering an environment that nurtures self-acceptance to promote healthier psychological development.
\\ \textit{Note:} correct because it accurately describes how Carl Rogers' concept of conditional positive regard influences an individual's self-concept and can lead to psychological maladjustment when external evaluations cause a disconnect between their true self and the self they portray.
\end{enumerate}

\vfill
\noindent\textit{End of Answer Key}
\end{document}